As demonstrated in the experimental results, the proposed computer vision pipeline proved robust across its interdependent modules. The boundary detection module successfully localized the playing surface in highly varied environments, overcoming significant differences in wood textures, surface reflections, camera angles, and ambient illumination. Furthermore, the ball tracking architecture reliably isolated and tracked the ball's kinematic trajectory, utilizing temporal interpolation and graph-based filtering to maintain continuity even during periods of fast motion or occlusion by the bowler. The optical flow-based spin estimation module effectively captured the rotational dynamics, yielding consistent rev rate profiles that successfully differentiated the unique bowling capacities and release styles of two distinct players. Ultimately, the integration of planar homography and geometric priors allowed for a metrically accurate and visually coherent 3D reconstruction of the ball's trajectory, successfully translating raw 2D pixel data into a reliable spatial model.

However, while the proposed lane detection pipeline demonstrates robustness across various illumination conditions, wood textures, and moderate camera angles, the reliance on rigid geometric priors inherently limits its flexibility. Empirical testing (e.g., Clip 8) revealed that severe deviations in camera composition, perspective, and optical zoom cause the foundational edge detection heuristics to fail. To address these limitations in highly unconstrained environments, future work could explore the integration of Deep Learning architectures. Developing a hybrid approach that combines the robustness of convolutional neural networks for initial boundary segmentation with our geometric algorithms for precise structural refinement presents a promising avenue for subsequent research.


Another limitation of the current 3D reconstruction stems from the planar assumption of the homography mapping, which mathematically restricts all tracked points to the floor plane (Z = 0). Consequently, if a bowler lofts the ball, the airborne pixel coordinates would be erroneously projected further down the lane, generating artificial spikes in the calculated spatial trajectory and velocity. To resolve this limitation, future research could evaluate the integration of physics-based motion priors to accurately estimate the ball's vertical height prior to surface contact.

Lastly, the system is constrained by the assumption of a static camera feed. Any camera movement would break the lane detection, trajectory detection, 3d reconstruction and spin detection modules, due to their underlying static hypothesis. To adapt the system for dynamic recording environments, subsequent research could explore frame-by-frame homography estimation or camera tracking algorithms to decouple sensor movements from the ball’s actual trajectory and spin metrics